\section{Workflow type: tracking mode and model}
\label{sec:workflow}

The \ui{WORKFLOW TYPE} section, below the image list (Figure~\ref{fig:workflow}), holds the two
choices that shape a run:
which image each frame is compared with (the \ui{Tracking Mode}) and which network compares them
(the \ui{Model}). It sits above \ui{REGION OF INTEREST} because the tracking mode decides which
frames are reference frames, and the hint of the region of interest follows it
(Section~\ref{sec:roi}). The card in the top-left corner of the canvas repeats both choices
(Section~\ref{sec:canvas}).

\screenshot[0.45\textwidth]{fig-workflow.png}{The \ui{WORKFLOW TYPE} section: the \ui{Tracking
Mode} and \ui{Model} lists and the \ui{Suggest a Model} button.}{fig:workflow}

\subsection{Tracking Mode}

A sequence of $n$ images is processed as $n-1$ image pairs, one for every frame after the first.
Each pair holds a \emph{reference frame} $r$ and a later frame $k$; the network measures the
displacement from $r$ to $k$. The tracking mode decides which reference frame each frame gets.
Every mode computes one pair per frame, and every mode reports its results in the coordinates of
frame~1. Table~\ref{tab:tracking-modes} summarises the choices.

\begin{table}[htbp]
  \centering
  \caption{What each frame $k$ is compared with. Frame numbers as in the image list.}
  \label{tab:tracking-modes}
  \begin{tabularx}{\textwidth}{>{\raggedright\arraybackslash}p{3.6cm}
                               >{\raggedright\arraybackslash}X
                               >{\raggedright\arraybackslash}p{3.4cm}}
    \toprule
    Tracking Mode & Frame $k$ is compared with & Reference frames \\
    \midrule
    \ui{Accumulative} & frame 1 & 1 \\
    \ui{Incremental}, \ui{Every Frame} & frame $k-1$ & every frame but the last \\
    \ui{Incremental}, \ui{Every N Frames} & frame $1 + N\lfloor (k-2)/N \rfloor$:
      frames 2 to $N+1$ with frame 1, frames $N+2$ to $2N+1$ with frame $N+1$, \ldots
      & $1, N+1, 2N+1, \ldots$ \\
    \ui{Incremental}, \ui{Custom Frames} & the last reference frame before $k$
      & 1 and the frames you list \\
    \bottomrule
  \end{tabularx}
\end{table}

\subsubsection*{Accumulative (default)}

Every frame is compared directly with frame~1. Nothing is composed, so an error in one frame
does not pass on to the next. This is the mode to use whenever the displacement from frame~1
stays within the reach of the model (Section~\ref{sec:workflow-model}) in every frame of the
sequence.

\subsubsection*{Incremental}

Each frame is compared with a reference frame nearer to it than frame~1, and the steps are
composed back to frame~1. Choose it when the motion from frame~1 grows beyond the reach of the model, as in a
large stretch or a large rotation. The price is that every step adds its own error to all the
frames measured after it.

Choosing \ui{Incremental} shows the \ui{Reference Update} list, which sets how often the
reference frame changes:
\begin{description}[leftmargin=1.2cm,style=nextline]
  \item[Every Frame (default)]
    Each frame is compared with the one before it. Each step is the motion between two
    consecutive images, the shortest step the sequence offers.
  \item[Every N Frames]
    Shows the \ui{Interval} box (2 to 100 frames, default 5). The reference frame changes every
    $N$ frames: with $N = 5$, frames 2 to 6 are compared with frame~1, frames 7 to 11 with
    frame~6, and so on. Fewer steps are composed than with \ui{Every Frame}, but each step is
    longer. Choose $N$ so that the motion over $N$ frames stays within the reach of the model.
  \item[Custom Frames]
    Shows the \ui{Reference Frames} field. Type the frame numbers, as in the image list,
    separated by commas (for example \file{5, 10, 20}); the setting applies when you press
    \kbd{Enter} or leave the field. Frame~1 is always the first reference frame, and each frame
    is compared with the last reference frame before it. Use it when the motion is uneven: put
    the reference frames closer together where the specimen moves fast. The last frame cannot be
    a reference frame, because no frame follows it, so with $n$ images the numbers must lie
    between 1 and $n-1$. A list that breaks this rule, or that is not made of whole numbers, is
    refused: the field gets a red border, its tooltip and a warning in the log say what is wrong,
    for example \emph{``Invalid reference frames 44: use frames 1 to 43, separated by commas (the
    last frame cannot be a reference).''}. While the field holds refused text, \ui{Run DIC
    Analysis} does not start: a \ui{Cannot start analysis} dialog repeats the reason. Reference
    frames
    that the loaded images cannot use, because a project or an earlier folder had more images,
    are dropped, and the log names them.
\end{description}

\subsubsection*{How the steps are composed}

For a pair $(r, k)$ the network returns the displacement $\mathbf{d}_{rk}$ from frame $r$ to
frame $k$, in the coordinates of frame $r$. A point $\mathbf{X}$ of the region of interest of
frame~1 sits at $\mathbf{X} + \mathbf{U}_r(\mathbf{X})$ in frame $r$, so its displacement to
frame $k$ is
\begin{equation*}
  \mathbf{U}_k(\mathbf{X}) = \mathbf{U}_r(\mathbf{X})
      + \mathbf{d}_{rk}\bigl(\mathbf{X} + \mathbf{U}_r(\mathbf{X})\bigr),
\end{equation*}
with $\mathbf{d}_{rk}$ interpolated bilinearly between pixels. In every mode the result is the
displacement of the points of frame~1 (Lagrangian coordinates): $+u$ to the right, $+v$
downward, in pixels.

Composition has three consequences worth knowing:
\begin{itemize}
  \item The region of interest of frame~1 is carried to each reference frame with the
        displacement measured so far, and the next pairs are solved over that carried region.
        A mask you give a reference frame replaces the carried region there
        (Section~\ref{sec:roi-later-frames}). At the start of a run, the \ui{LOG} lists the
        reference frames that carry the region of frame~1: \emph{Reference frames without a
        Region of Interest of their own carry the one of frame 1 forward}, followed by their
        numbers.
  \item A point that is lost in a reference frame (it left the image, or a mask excludes it)
        stays empty in every frame measured from that reference frame.
  \item If a reference frame has no result, the frames compared with it are skipped, and the
        log says so.
\end{itemize}

\begin{notebox}
After a run, RAFTcorr watches the residual (the grey-level mismatch left after the
displacement maps one image onto the other) and warns when it jumps from one frame on. In
accumulative tracking the warning suggests incremental tracking or a model with a larger $s$;
in incremental tracking it suggests more reference frames or a model with a larger $s$
(Section~\ref{sec:tracking-quality}).
\end{notebox}

\subsection{Model}
\label{sec:workflow-model}

The \ui{Model} list holds the four models of Table~\ref{tab:models}, each with its reach in
parentheses, as \ui{RAFTcorr-s4 (128 px)}. The reach is the nominal lookup radius of the model:
its coarsest correlation level looks $32\,s$~px away from each point. A smaller $s$ resolves finer
detail, such as a strain concentration, and a larger $s$ reaches further, so pick the finest model
whose reach covers the largest displacement of any pair. In accumulative tracking that is the
largest displacement from frame~1; in incremental tracking it is the largest step between a frame
and its reference frame.

Hovering over an entry shows its description and reach, and \ui{Weights not installed.} when the
model file is not on the computer; the list checks this each time it opens, so a model you have
just downloaded or imported shows as installed. A model is downloaded the first time a run or a survey
uses it; \menu{Help}{Models…} shows which models are installed and imports model files on a
computer without an internet connection (Section~\ref{sec:installing-models}).

\subsection{Suggest a Model}
\label{sec:suggest}

\ui{Suggest a Model} measures how far the specimen moves and selects a model for you. Use it
before the first run, in particular when you do not know the size of the motion. It needs at
least two images and does not start while a run is going. It works in these steps:
\begin{enumerate}
  \item It loads RAFTcorr-s8, the model with the largest reach (downloading it on first use),
        with the current \ui{Device} and \ui{Iterations} (Section~\ref{sec:parameters}). A
        progress dialog, \ui{Measuring the motion with RAFTcorr-s8}, has a \ui{Cancel} button.
  \item It compares up to three pairs of frames, chosen by the tracking mode. In
        \ui{Accumulative} tracking these are frame~1 against the frame a quarter of the way
        through the sequence, the middle frame and the last frame (for 44 images, frames 11, 22
        and 44). In \ui{Incremental} tracking it takes the reference frames the run would use and,
        for each, the frame farthest from it that is compared with it, where the step is
        largest; of these pairs it surveys the last, the middle one and the one a quarter of the
        way through. The progress dialog names each pair, for example \ui{Survey: frame 44 against
        frame 43}.
  \item Over the region of interest of frame~1 (the whole image if none is drawn yet), it takes
        the 99th percentile of the displacement magnitude $\sqrt{u^2 + v^2}$ of each sampled
        pair. The 99th percentile ignores the largest 1\,\% of the values, so a few stray points
        do not decide the choice. The largest of the three values is the measured displacement
        $d$.
  \item It selects the finest model whose reach is at least $1.3\,d$: RAFTcorr-s1 for $d$ up to
        about 24.6~px, -s2 up to about 49.2~px, -s4 up to about 98.5~px, and -s8 above.
\end{enumerate}

The \ui{Model} list switches to the suggestion, and the same message appears in the \ui{LOG}
and in a card under the button (the parts in angle brackets are the measured values):
\begin{lstlisting}
Suggested: RAFTcorr-<s>. Largest displacement of the region of interest about <d> px (99th percentile over frames <a, b, c>); RAFTcorr-<s> reaches <R> px.
\end{lstlisting}
In incremental tracking the message lists the surveyed pairs, each as reference frame
$\rightarrow$ frame:
\begin{lstlisting}[literate={→}{{$\rightarrow$}}1]
Suggested: RAFTcorr-<s>. Largest displacement of the region of interest between a frame and its reference frame about <d> px (99th percentile over the pairs reference → frame <r1 → k1, r2 → k2, r3 → k3>); RAFTcorr-<s> reaches <R> px.
\end{lstlisting}
When $1.3\,d$ exceeds even the reach of RAFTcorr-s8 (above about 197~px), the suggestion is still
RAFTcorr-s8, and the message adds a sentence that starts with \ui{This exceeds even the reach of
RAFTcorr-s8:} and says what to change: \ui{use incremental tracking.} in accumulative tracking,
\ui{update the reference frame more often.} with \ui{Every N Frames} or \ui{Custom Frames}, and
\ui{the motion between consecutive frames is too large to track.} with \ui{Every Frame}. The card
stays until you open other images, and a project file keeps it.

\begin{tipbox}
Draw the region of interest before you click \ui{Suggest a Model}. The survey then measures the
motion of the specimen, not of the background around it.
\end{tipbox}

\begin{notebox}
The survey looks at three pairs only. If the motion peaks between the sampled frames and comes
back, or one step of an incremental run is much longer than the surveyed ones, the survey does
not see it. Choose the tracking mode before you click \ui{Suggest a Model}: a suggestion made
for accumulative tracking does not fit an incremental run, and the other way round.
\end{notebox}

\subsection{Which combination to choose}

\begin{center}
\begin{tabularx}{\textwidth}{>{\raggedright\arraybackslash}X
                             >{\raggedright\arraybackslash}p{4.2cm}
                             >{\raggedright\arraybackslash}p{3.6cm}}
  \toprule
  Situation & Tracking Mode & Model \\
  \midrule
  The displacement from frame~1 stays within the reach of a model in every frame
    & \ui{Accumulative} & the finest model that reaches it (\ui{Suggest a Model}) \\
  The displacement from frame~1 exceeds the reach of RAFTcorr-s8
    & \ui{Incremental}, \ui{Every Frame} & the finest model that reaches the step between
      two images \\
  Many images, each step small, large total motion
    & \ui{Incremental}, \ui{Every N Frames} & a model whose reach covers the motion over
      $N$ frames \\
  Uneven motion (slow, then fast)
    & \ui{Incremental}, \ui{Custom Frames} & a model whose reach covers the longest step \\
  \bottomrule
\end{tabularx}
\end{center}
